% !TeX program = pdflatex
% !TeX encoding = utf8
% !TeX spellcheck = uk_UA


\documentclass{article}
\usepackage[utf8]{inputenc}
\usepackage[ukrainian]{babel}
\usepackage{amsmath}
\usepackage{amssymb}
\usepackage{mhchem}

\begin{document}

\section*{88. Віріальне розкладання та температура Бойля}

Рівняння стану ідеального газу, з яким ми оперували вивчаючи термодинамічний метод, очевидно неповно
описує поведінку реальних газів у всьому діапазоні температур і тисків. Як мінімум, це рівняння в
принципі не здатне описати стан речовини при низьких температурах, коли будь-яка речовина має
перебувати в сконденсованому стані. Звернемося до дослідних даних. Як уже зазначалося в першій лекції,
принципово, поведінку будь-якого газу можна описати за допомогою \emph{віріального розвинення},
запропонованого Камерлінг-Оннесом:

\begin{equation}
pV = RT + Bp + Cp^{2} + \ldots \quad \propto \quad pV = RT + \frac{B'}{V} + \frac{C'}{V^{2}} + \ldots
\label{eq:88.1}
\end{equation}

де коефіцієнти $A=RT, B, C, B', C'$ -- віріальні коефіцієнти. Із цих коефіцієнтів, $B$ дає більш
істотну поправку у відмінності рівняння стану реальних газів від ідеальних. Так, для азоту при
$50^\circ\mathrm{C}$, виражаючи $P$ в мм.рт.ст., і $V$ в літрах маємо: $A = 1.0$; $B = -0.61 \cdot
10^{-3}$; $C = 5.4 \cdot 10^{-6}$. Відклавши на графіку $B = f(T)$ експериментальні значення другого
віріального коефіцієнту $B$, виявимо, що залежність ця має однаковий вигляд для всіх газів: різке
зростання при відносно низьких температурах і подальша стабілізація і іноді ледь помітне спадання.

Температура, при якій $B = 0$, отримала назву \emph{температури Бойля}. При цій температурі, газ, з
точністю до наступного коефіцієнта віріального розвинення, задовольняє, подібно до ідеального газу
закону Бойля-Маріотта. Крива, для якої $\left(\frac{\partial(PV)}{\partial P}\right)_T = 0$, отримала
назву \emph{кривої Бойля}. В околицях цієї кривої газ поводиться подібно до ідеального. На залежності
$PV = f(P)$ крива утворюється геометричним місцем точок мінімумів залежності $PV = f(P)$ для різних
температур.

\section*{89. Ізотерми Ендрюса та критичні параметри}

Переконлива відмінність поведінки реальних газів від ідеальних була встановлена Томасом Ендрюсом, при
побудові ізотерм \ce{CO_2} (\emph{ізотерми Ендрюса}). Згодом подібна поведінка була досліджена для
всіх газів. Якщо вуглекислоту стискати при $t = 21.5^\circ\mathrm{C}$, то після деякого ступеня
стиснення $P$, газ починає стискатися. Тиск при цьому залишається постійним, аж до моменту повного
зрідження газу. Явища ці спостерігаються аж до температури $30.9^\circ\mathrm{C}$ і вже при
$31^\circ\mathrm{C}$, вуглекислоту не вдається сконденсувати ні при яких тисках.

У зв'язку з цим температуру $30.9^\circ\mathrm{C}$ називають \emph{критичною температурою} для
вуглекислоти, а ізотерму при $t = 30.9^\circ\mathrm{C}$ -- \emph{критичною ізотермою}. Всі ізотерми,
що відповідають високим температурам, увігнуті. На критичній ізотермі існує точка перегину. Значення
$P$ і $V$, які їй відповідають, отримали назву критичних: $P_{\text{кр}}$, $V_{\text{кр}}$. Всі гази
поводять себе в подібних дослідах однаково -- значення ж $P_{\text{кр}}$, $V_{\text{кр}}$ і
$T_{\text{кр}}$ різняться в широких межах.

На останок зазначимо, що крива (позначена пунктиром), яка обмежує двофазну область, вершиною якої і є
критична точка, називається \emph{кривою лабільності}.

\section*{90. Рівняння Ван-дер-Ваальса}

Оскільки симетрія рідкого і газоподібного стану неперервна, то можна сподіватися отримати загальне
рівняння стану, яке може застосовуватися до обох станів речовини. Зауважимо, що цього не можна зробити
для рідини і твердого тіла внаслідок істотно різних симетрій, властивих рідині та твердому тілу. Ян
Дидерік ван дер Ваальс в своїй роботі «\emph{Неперервність газоподібних і рідких станів}» в 1873 р.,
вперше не тільки запропонував рівняння, що задовільно описує емпіричні дані, а й дав фізичне пояснення
поправочних коефіцієнтів, які входять до рівняння В.д.В:

\begin{equation}
\left( p + \frac{a}{v^{2}} \right)(v - b) = RT\quad
v = \frac{V}{\nu}
\label{eq:90.1}
\end{equation}

де $v$ -- об'єм, який припадає на один моль газу і для газу, який знаходиться при нормальних умовах,
рівний
$v = \frac{22.4\ \text{л}}{6.02 \cdot 10^{23}} = 3.72 \cdot 10^{-23}$~см$^{3}$.
Звідси знаходимо, що середня відстань між молекулами $\langle r \rangle \approx 33\cdot10^{-7}
\text{см}$. Зроблена оцінка дозволяє оцінити ступінь наближення, припустимого в моделі ідеальних
газів, де
розміри молекул не беруться до уваги, в той час як «розміри» реальних молекул, а точніше характерні
масштаби взаємодії молекул, згідно квантово-механічному підходу, складають $0.5\div5\cdot10^{-8}$~см.

Питання про характер або навіть знак взаємодії послідовно вирішується методами квантової механіки.
Однак, деякі висновки можна зробити аналізуючи результати дослідів з реальними газами. Зауважимо, що
знак похідної $\left(\frac{\partial(PV)}{\partial P}\right)_T$, що можна знайти із ізотерм $PV =
f(P)$, вказує на відносну стисливість газу:

\begin{equation}
\left( \frac{\partial(pV)}{\partial p} \right)_{T} = V + p \cdot \left( \frac{\partial V}{\partial P}
\right)_{T} = V(1 - p \cdot \beta) = V\left( 1 - \frac{\beta}{\beta_{ід}} \right)
\label{eq:90.2}
\end{equation}

Як видно з графіка $PV = f(P)$ при малих об'ємах та високих температурах ізотермічна стисливість
реального газу менша, ніж у ідеального $\left(\frac{\partial(PV)}{\partial P}\right)_T > 0$. Тобто при
малих відстанях переважають сили відштовхування, що діють між молекулами і ускладнюють стиснення газу.
При високих температурах збільшується частота зіткнень між молекулами. Завдяки цьому молекули досить
часто опиняються близько одна від одної, коли відстань між ними мала і діють сили відштовхування.

Область на графіку $PV = f(P)$ зліва від кривої Бойля відповідає станам, де ізотермічна стисливість
реального газу більша, ніж у ідеального. У цій області домінують сили притяжіння між молекулами, що
сприяють стисненню газу. Ці сили отримали назву \emph{вандервальсових}. Вони залежать від типу
взаємодіючих молекул (полярні, неполярні), хоча в кінцевому рахунку їх природа це диполь-дипольна
взаємодія. При дуже малих відстанях між молекулами, електронні оболонки взаємодіючих молекул
перекриваються і, як показують квантово-механічні розрахунки, сили притяжіння переходять в сили
відштовхування, які різко зростають зі зменшенням відстані. Схематично потенціал взаємодії $u(r)$ між
молекулами виглядає як на рисунку. Найчастіше його апроксимують \emph{потенціалом Леннарда-Джонса}:

\begin{equation}
u(r) = \frac{\alpha}{r^{12}} - \frac{\beta}{r^{6}}
\label{eq:90.3}
\end{equation}

\section*{91. Модель Ван-дер-Ваальса}

Перший член в \eqref{eq:90.3} настільки швидко зростає, що його дію можна замінити дією нескінченно
високої стінки в точці $\sigma = R$ — радіусу молекули. Так і роблять у теорії ван дер Ваальсу,
апроксимуючи реально діючий потенціал відштовхування моделлю абсолютно пружних куль певного радіуса
$R$.

Виділимо деяку область в об'ємі газу, в яку входять лише дві молекули газу. При зіткненні центри
молекул не можуть зблизитися ближче ніж на $D = 2R$. Це означає, що в моделі Ван-дер-Ваальсу, ми,
розглядаючи взаємодію молекул при зіткненнях й дотепер можемо їх вважати точковими, проте доступний їм
об'єм зменшується. Так, в розглянутому випадку взаємодії двох молекул, недоступний об'єм становить
$\frac{4}{3}\pi D^3 = 4\cdot 2V_{м}$ — учетверенний об'єм двох молекул.

У сусідній області, яка обмежує дві інші молекули, недоступний їм об'єм складе ту ж величину. Всього,
об'єм недоступний чотирьом молекулам складе $4.4V_{м}$. Розбиваючи область, зайняту газом на області
взаємодії, отримаємо, що загальний недоступний $N_A$ молекулам газу об'єм становить $b = 4N_A V_{м}$.

Зауважимо, що в наведеному розгляді не враховуються випадки потрійних, четверних та інших зіткнень,
частка яких, зрозуміло, мізерно мала. Тобто врахування сил відштовхування можна зробити, модифікуючи
рівняння Клайперона-Мендєлєєва — зменшивши об'єм газу на величину об'єму недоступного молекулам:

\begin{equation}
P(V - \nu b) = \nu RT
\label{eq:91.1}
\end{equation}

Для врахування сил тяжіння зауважимо, що молекули пристінкового шару і молекули всередині газу
знаходяться в різних умовах. На молекулу 1 на рисунку сили тяжіння діють з усіх боків, в той час як на
молекулу 2 поблизу стінки діють сили з боку молекул газу, рівнодіюча яких прагне «втягнути» молекулу в
газ. Тиск, який ми вимірюємо, створюється молекулами пристінкового шару і він завжди менший за
внутрішній. Тобто $P = P_{\text{in}} - \Delta P$.

Ми як і раніше можемо використовувати рівняння стану ідеального газу з поправкою на недоступний об'єм
\eqref{eq:91.1}, проте замість вимірюваного тиску $P$ слід використовувати $P_{\text{in}}$. Можна
заперечити, що на молекули пристінкового шару діють також сили з боку стінки. Однак, на відміну від
сил тяжіння з боку молекул газу, їх дія визначає параметри взаємодії молекул стінки і молекул газу
(пружнє, непружнє, є залипання чи ні і т.д.). Якщо ж вже приймається, що ця взаємодія пружна, то вид
її може вплинути тільки на характер руху молекул поблизу стінки. Для створюваного тиску важливо
тільки, що нормальна складова імпульсу після зіткнення змінює знак.

Звернемося тепер до оцінки $\Delta P$. Величина ця пропорційна силі, яка втягує молекулу всередину, а
також кількості молекул поблизу стінки. Обидві ці величини пропорційні концентрації молекул (або
щільності газу). Звідси $\Delta P = \frac{a\nu^2}{V^2}$ і в результаті рівняння набуває вигляду:

\begin{equation}
\left( P + \frac{a\nu^{2}}{V^{2}} \right)(V - \nu b) = \nu RT
\label{eq:91.2}
\end{equation}

\section*{92. Інші рівняння стану реального газу}

Слід зазначити, що модель газу Ван-дер-Ваальса всього лише модель, яка більш-менш точно описує стан
реального газу. Існують і інші рівняння стану реального газу. Серед інших відзначимо \emph{рівняння
Дітерічі}:

\begin{equation}
P(V - b) = RT\exp\left\{ - \frac{a}{RTV} \right\}
\label{eq:92.1}
\end{equation}

яке більш точно описує поведінку газу при зменшених тисках, але не дає переваг при високих. Якщо $b
\ll V$ і $a \ll RTV$, то рівняння Дітерічі переходить в рівняння Ван-дер-Ваальса.

Ще одне рівняння — \emph{рівняння Бертло}:

\begin{equation}
\left( P + \frac{a}{TV^{2}} \right)(V - b) = RT
\label{eq:92.2}
\end{equation}

також більш точно описує стан газу при низьких тисках, але не дає переваг поблизу критичної точки.

Ще більш точним є \emph{рівняння Клаузіуса}:

\begin{equation}
\left( P + \frac{a}{T(V + c)^{2}} \right)(V - b) = RT
\label{eq:92.3}
\end{equation}

за рахунок того, що вводиться ще одна (третя) емпірична поправка $c$.

\section*{93. Аналіз рівняння Ван-дер-Ваальса}

Проаналізуємо рівняння Ван-дер-Ваальса і встановимо, наскільки точно воно описує реальні гази.
Перепишемо (91.2) як:

\begin{equation}
PV^{3} - (RT + Pb)V^{2} + aV - ab = 0
\label{eq:93.1}
\end{equation}

Це парабола третього ступеня, в яку $P$ і $T$ входять у вигляді параметрів. При великих $T$ членами
$Pb$, $aV$, $ab$ можна знехтувати і рівняння (1) описує гіперболу -- ізотерми ідеального газу. У
загальному випадку рівняння (1) має або три дійсних кореня, або один. Кожному кореню відповідає точка
на площині $(P, V)$, де ізобара перетинає кубічну параболу. Існує лише одна точка на площині $(P, V)$
де всі три корені збігаються. Очевидно, що (1) за теоремою Вієта тоді можна подати так:

\begin{equation}
P_{к}(V - V_{к})^{3} = 0
\label{eq:93.2}
\end{equation}

Розкриваючи (2) і порівнюючи з (1) знайдемо значення, які набувають $P$, $V$, $Т$ в цій особливій
точці:

\begin{align}
V_{K} &= 3b \label{eq:93.3} \\
P_{K} &= \frac{a}{27b^{2}} \label{eq:93.4} \\
T_{K} &= \frac{8a}{27Rb} \label{eq:93.5}
\end{align}

Зауважимо, що тільки в цій точці виконуються умови:

\begin{equation}
\left( \frac{\partial P}{\partial V} \right)_{T} = \left( \frac{\partial^{2}P}{\partial V^{2}}
\right)_{T} = 0
\label{eq:93.6}
\end{equation}

Тобто зазначена точка є ніщо інше, як критична точка, яка є характерною точкою на ізотермах Ендрюса.
Величину:

\begin{equation}
K_{K} = \frac{RT_{K}}{P_{K}V_{K}} = \frac{8}{3} = 2.67
\label{eq:93.7}
\end{equation}

називають \emph{критичним коефіцієнтом}. Для реальних газів ця величина змінюється від 3.5 до 4.5.
Збіг можна визнати в міру задовільним.

Подивимося, яким чином співвідносяться теоретично передбачена і експериментально визначена температура
Бойля. Для цього запишемо (1) у вигляді:

\begin{equation}
pV = \frac{RT}{(1 + \frac{a}{pV^{2}})(1 - \frac{b}{V})}
\label{eq:93.8}
\end{equation}

Розкладаючи в ряд Тейлора і нехтуючи членами другого порядку малості, а також вважаючи, там де
необхідно, що $PV \approx RT$ отримаємо з (6):

\begin{equation}
PV = RT\left( 1 - \frac{a}{PV^{2}} + \frac{a}{PV^{3}}\frac{b}{V} \right) = RT - \left( \frac{a}{RT} -
b \right)P - \frac{ab}{RT}P^{2} = RT + BP + CP^{2} + \ldots
\label{eq:93.9}
\end{equation}

Звідки знаходимо що $T_{B} = a / Rb$, а відношення $T_{B} / T_{K} = 27/8$:

\begin{equation}
T_{B} = \frac{27}{8}T_{K}
\label{eq:93.10}
\end{equation}

Для реальних газів це відношення змінюється від $3.65$ для гелію до $2.56$ для азоту. Тобто тут збіг
можна
вважати лише задовільним. Проте рівняння Ван-дер-Ваальсу якісно вірно передає поведінку речовини аж до
його зрідження, а з деякими застереженнями, які ми наведемо далі, і в двофазній області.


\section*{94. Закон відповідних станів}

Наявність особливої точки на площині $(P, V)$ дозволяє ввести природні одиниці виміру тиску — $P_{K}$,
об'єму — $V_{K}$ і температури — $T_{K}$. Підставляючи їх значення (93.3) в рівняння Ван-дер-Ваальса
(91.2), отримаємо для одного моля речовини:

\begin{equation}
\left( \pi + \frac{3}{\varphi^{2}} \right)\left( \varphi - \frac{1}{3} \right) = \frac{8}{3}\tau
\label{eq:94.1}
\end{equation}

де $\pi = P / P_{K}$; $\varphi = V / V_{K}$; $\tau = T / T_{K}$.

У це рівняння не входять жодні сталі, які характеризують речовину (окрім власне критичних параметрів).
Отже, наведене рівняння вірно для всіх речовин, якщо реальні тиск, об'єм і температуру виразити в їх
критичних значеннях. Отриманий висновок є більш загальним, ніж саме рівняння стану Ван-дер-Ваальса.

Справді, будь-які з теоретично або емпірично отриманих функцій стану, що залежать від трьох емпіричних
параметрів, можуть бути подані в безрозмірному вигляді. Подаючи емпіричні результати в одиницях
$P_{K}$, $T_{K}$, $V_{K}$, ми отримаємо можливість порівнювати експеримент і теорію або різні теорії.

Положення про те, що рівняння стану, записане в безрозмірних одиницях $\pi$, $\varphi$, $\tau$, має
бути однаковим для всіх речовин, називають \emph{законом відповідних станів}: якщо для різних речовин
дві з величин збігаються, то буде збігатися і третя. Цей закон є одним із прикладів \emph{методу
подібності}, часто використовуваного в фізиці.

Зокрема, подавши криві $PV = f(P)$ в одиницях $PV/RT$ та $P/P_{K}$, зручно порівнювати криві, отримані
для різних газів. Також, наприклад, теплоємність твердого тіла, виражена в одиницях $R$, або
намагніченість феромагнетиків, віднесена до намагніченості при $T = 0\,\text{K}$, має однакову
температурну залежність, якщо температура вимірюється в одиницях температури Дебая ($T_{D}$) в першому
випадку і в одиницях температури Кюрі в другому випадку.


\section*{95. Термодинамічні властивості газу Ван-дер-Ваальса}

Розглянемо тепер деякі термічні та калоричні властивості газу Ван-дер-Ваальса. Коефіцієнт лінійного
розширення дорівнює:

\begin{equation}
\alpha = \frac{1}{V}\left( \frac{\partial V}{\partial T} \right)_{P} = \frac{(V - b)}{VT -
\frac{2a}{R}\left( \frac{V - b}{V} \right)^{2}}
\label{eq:95.1}
\end{equation}

Для ідеального газу ($a = b = 0$) отримуємо $\alpha = 1/T$. Порівняємо коефіцієнт $\alpha$ газу
Ван-дер-Ваальса з коефіцієнтом $\alpha$ ідеального газу. Вважаючи, що $V \gg b$ і $RTV \gg a$,
розкладемо (\ref{eq:95.1}) в ряд:

\begin{equation}
\alpha - \frac{1}{T} = \frac{\left( \frac{2a}{RT} - b \right)}{VT} = \frac{b}{V}\frac{1}{T}\left(
\frac{2T_{B}}{T} - 1 \right)
\label{eq:95.2}
\end{equation}

Отже, при температурі, яка дорівнює подвійній температурі Бойля, коефіцієнти зрівнюються.

Використовуючи основний наслідок другого закону термодинаміки, знаходимо:

\begin{equation}
\left( \frac{\partial U}{\partial V} \right)_{T} = T\left( \frac{\partial P}{\partial T} \right)_{V} -
P = \frac{RT}{V - b} - \frac{RT}{V - b} + \frac{a}{V^{2}} = \frac{a}{V^{2}}
\label{eq:95.3}
\end{equation}

Цей результат є очікуваним, оскільки в моделі враховано сили відштовхування та притягання, які
по-різному залежать від відстані. Інтегруючи (\ref{eq:95.3}), отримаємо:

\begin{equation}
U = -\frac{a}{V} + \phi(T)
\label{eq:95.4}
\end{equation}

З іншого боку, $\left( \frac{\partial U}{\partial T} \right)_{V} = C_{V} = \frac{\partial
\phi}{\partial T}$. Таким чином, вираз для внутрішньої енергії газу Ван-дер-Ваальса має вигляд:

\begin{equation}
U = C_{V}T - \frac{a}{V} + U_{0}
\label{eq:95.5}
\end{equation}

Отже, для газу Ван-дер-Ваальса внутрішня енергія є функцією температури та об'єму. Зауважимо, однак,
що, як і для ідеального газу, теплоємність $C_{V}$ не залежить від об'єму:

\begin{equation}
\frac{\partial C_{V}}{\partial V} = \frac{\partial^{2}U}{\partial V \partial T} =
\frac{\partial^{2}U}{\partial T \partial V} = \frac{\partial}{\partial T}\left( \frac{a}{V^{2}}
\right) = 0
\label{eq:95.6}
\end{equation}

Для обчислення ентропії підставимо вираз для внутрішньої енергії (\ref{eq:95.5}) у перший закон
термодинаміки:

\begin{equation}
dS = \frac{C_{V}dT}{T} + \left[ P + \frac{a}{V^{2}} \right]\frac{dV}{T} = \frac{C_{V}dT}{T} +
\frac{R}{V - b}dV
\label{eq:95.7}
\end{equation}

Інтегруючи (\ref{eq:95.7}), знайдемо:

\begin{equation}
S = C_{V}\ln T + R \ln (V - b) + \text{const}
\label{eq:95.8}
\end{equation}

Нарешті, знайдемо різницю $C_{P} - C_{V}$, скориставшись отриманим раніше виразом для різниці
теплоємностей:

\begin{equation}
C_{P} - C_{V} = T\left( \frac{\partial P}{\partial T} \right)_{V}\left( \frac{\partial V}{\partial T}
\right)_{P} = \frac{RT(V - b)}{(V - b)\left( 1 - \frac{2a}{RT}\frac{(V - b)^{2}}{V^{3}} \right)}
\approx \frac{R}{1 - \frac{2a}{RTV}}
\label{eq:95.9}
\end{equation}

\section*{96. Фазові переходи та правило Максвела}

Розглянемо тепер хід ізотерм у області, де вони мають осциляції. Очевидно, що ділянка \textit{de} не
може бути реалізована, оскільки на ній $\left(\frac{\partial P}{\partial V}\right)_T > 0$, що є
неможливим для термодинамічно стійкої системи. Експериментальні результати свідчать, що реалізується
ізобаричні-ізотермічні ділянки \textit{ac}. Це область двофазного стану.

Умова співіснування газоподібної та рідкої фаз визначається рівністю питомих термодинамічних
потенціалів Гіббса в обох фазах. Розглянемо ізотерму-ізобару, де $g_a = g_c$. Тоді:

\begin{equation}
f_a + PV_a = f_c + PV_c
\label{eq:96.1}
\end{equation}

Оскільки:

\begin{equation}
f_a - f_c = -\int_{a}^{c} \left( \frac{\partial f}{\partial V} \right)_T dV = \int_{a}^{c} P dV
\label{eq:96.2}
\end{equation}

отримаємо:

\begin{equation}
P(V_c - V_a) = \int_{a}^{c} P dV
\label{eq:96.3}
\end{equation}

Отримане співвідношення показує, як слід проводити ізотерму-ізобару \textit{ac}. Ліва частина рівняння
(\ref{eq:96.3}) дорівнює площі під ізобарою \textit{ac}, а права частина — площі під ізотермою
\textit{adbec}. Отже, ізотерму \textit{ac} слід проводити так, щоб площа \textit{adb} дорівнювала
площі \textit{bec}. Це правило називається \emph{правилом Максвела}.

Зауважимо, що при побудові не використовувався явний вигляд рівняння Ван-дер-Ваальса. Наведена
побудова справедлива для будь-якого рівняння стану. Також слід зазначити, що ділянки \textit{ad} та
\textit{ec}, хоча механічно стійкі, на практиці не реалізуються або потребують особливих умов для
реалізації. Ділянка ізотерми \textit{ec} називається областю \emph{перегрітої пари}, а \textit{ad} —
областю \emph{переохолодженої рідини}.

\section*{97. Правило важеля}

Розглянемо фізичний зміст різних точок ізотерми \textit{ac}. У різних точках відрізка присутні різні
частки рідкої та газоподібної фаз. Точка \textit{c} описує чистий насичений пар з молярним об'ємом
$v_c$, а точка \textit{a} — чисту рідину з молярним об'ємом $v_a$.

Якщо позначити відносну кількість рідини через $x$, то:
\begin{itemize}
\item У точці \textit{c}: $x = 0$
\item У точці \textit{a}: $x = 1$
\end{itemize}

У довільній точці \textit{P} ізотерми \textit{ac} маємо:

\begin{equation}
v_P = x v_a + (1 - x) v_c
\label{eq:97.1}
\end{equation}

Звідси:

\begin{align}
x &= \frac{v_c - v_P}{v_c - v_a} = \frac{pc}{ac} \label{eq:97.2} \\
1 - x &= \frac{v_P - v_a}{v_c - v_a} = \frac{ap}{ac} \label{eq:97.3} \\
\frac{x}{1 - x} &= \frac{pc}{ap} \label{eq:97.4}
\end{align}

Таким чином, відносні частки рідкої та газоподібної фаз відносяться як відповідні відрізки. Це правило
визначення відносних часток називається \emph{правилом важеля}.


\section*{98. Визначення критичних параметрів}

Проводячи за правилом Максвела ізобари при різних температурах, можна для кожної ізотерми встановити
питомі об'єми рідини та газу, тим самим визначити область лабільності — область двофазних станів.
Зліва від неї речовина знаходиться лише в рідкому стані, справа — лише в газоподібному. Вершиною цієї
області є критична точка. Очевидно, цей метод застосовується до будь-якої речовини незалежно від того,
наскільки добре вона описується рівнянням Ван-дер-Ваальса.

Практичне визначення параметрів критичної точки здійснюється наступним чином. Ампулу, заповнену
частково певною кількістю речовини, нагрівають і визначають температуру, при якій зникає межа між
рідиною та парою. Це і буде критична температура. На перший погляд може здатися, що кількість речовини
в ампулі має бути чітко визначеною, проте це не так. Наприклад, для спостереження ефекту зникнення
мениска у вуглекислоті допустима зміна густини становить 0,735–1,269 від критичної.

Справа в тому, що поблизу критичної точки стисливість речовини прямує до нескінченності. Навіть
незначна зміна тиску, обумовлена дією верхніх шарів речовини на нижні, призводить до того, що густина
нижніх шарів може бути на десятки відсотків вищою за густину верхніх. Отже, завжди знайдеться шар
речовини, густина якого точно відповідає критичній.

Після визначення критичної температури, досліджуючи залежність $P = f(V)$ на критичній ізотермі, можна
знайти $P_{\text{кр}}$. Значення $V_{\text{кр}}$ або критичної густини найточніше визначається
\emph{методом прямолінійного діаметра}. У цьому методі густини рідини та газу вимірюють якомога ближче
до критичної точки. Експериментальні дані, нанесені на графік $\rho = \rho(T)$, утворюють викривлену
параболу. Значення $(\rho_{\text{р}} + \rho_{\text{г}})/2$, екстрапольоване до критичної температури,
дає критичну густину.

Крім властивості нескінченно високої стисливості речовини в критичній точці, існує низка інших цікавих
властивостей:
\begin{itemize}
\item Термодинамічна рівновага в цій точці встановлюється надзвичайно повільно
\item Процес зникнення мениска може бути досить тривалим
\item Явище критичної опалесценції проявляється у значному розсіюванні світла
\end{itemize}

Ці ефекти обумовлені гігантськими флуктуаціями, зокрема густини. Цими ж причинами пояснюється
прямування теплоємності $C_V$ до нескінченності при $T \to T_{\text{кр}}$. Теплоємність $C_P$ також
прямує до нескінченності, оскільки:

\begin{equation}
C_P - C_V = -T\left(\frac{\partial P}{\partial T}\right)_V^2 \left(\frac{\partial V}{\partial
P}\right)_T \to \infty
\end{equation}

внаслідок того, що $\left(\frac{\partial P}{\partial V}\right)_T \to 0$ при $T = T_{\text{кр}}$.

\section*{99. Ефект Джоуля-Томсона}

Вже у відомих дослідах Джоуля було встановлено, що при протіканні газу через пористу перегородку
температура більшості газів знижується, хоча й на дуже малу величину — тим меншу, чим більш газ
наближений до ідеального. Як було встановлено, цей процес є ізоентальпійним, отже:

\begin{equation}
\left(\frac{\partial T}{\partial P}\right)_H = \frac{V}{C_P}(1 - \alpha T)
\label{eq:99.1}
\end{equation}

Використовуючи вираз (95.2), знайдемо, що для газу Ван-дер-Ваальса коефіцієнт диференціального ефекту
Джоуля-Томсона дорівнює:

\begin{equation}
\mu = \left(\frac{\partial T}{\partial P}\right)_H = \frac{\left(\frac{2a}{RT} - b\right)}{C_P} =
\frac{b}{C_P}\left(\frac{2T_B}{T} - 1\right)
\label{eq:99.2}
\end{equation}

Звідси випливає, що газ охолоджується, якщо $T < 2T_B$, і нагрівається, якщо $T > 2T_B$. При кімнатних
температурах лише два гази (водень і гелій) нагріваються при дроселюванні. Температура $T_i = 2T_B$
називається \emph{температурою інверсії диференціального ефекту Джоуля-Томсона}.

Як випливає з (\ref{eq:99.1}), у загальному випадку $\left(\frac{\partial T}{\partial P}\right)_H =
0$, якщо $\alpha = 1/T$. Взявши за основу вираз для коефіцієнта об'ємного розширення газу
Ван-дер-Ваальса (95.1), знайдемо криву інверсії:

\begin{equation}
\frac{2a}{RT}\left(\frac{v - b}{v}\right)^2 = b
\label{eq:99.3}
\end{equation}

або в наведених одиницях:

\begin{equation}
\pi = 24\sqrt{3\tau} - 12\tau - 27
\label{eq:99.4}
\end{equation}

Отримана крива визначає область, де відбувається охолодження газу, та область, де спостерігається
нагрівання. Ці результати особливо важливі з технічної точки зору, оскільки дозволяють визначити
умови, за яких можливе зрідження газу методом Джоуля-Томсона.

Для визначення величини зниження температури газу при значних перепадах тиску розглядають
\emph{інтегральний ефект Джоуля-Томсона}. Величина $\Delta T$ визначається інтегруванням:

\begin{equation}
T_2 - T_1 = \int_{P_1}^{P_2} \left(\frac{\partial T}{\partial P}\right)_H dP = \int_{P_1}^{P_2}
\frac{V(\alpha T - 1)}{C_P} dP
\label{eq:99.5}
\end{equation}

У технічних пристроях тиск $P_2$ зазвичай має порядок атмосферного, $T_1$ — температуру попереднього
охолодження. Для знаходження умов максимального зниження температури необхідно виконати:

\begin{equation}
\left(\frac{\partial T}{\partial P}\right)_H = 0
\label{eq:99.6}
\end{equation}

Це означає, що точка $(P,T)$ лежить на кривій інверсії. Відповідно до цього принципу конструюються
холодильні машини. Наприклад, при зріджуванні водню попереднє охолодження досягається кипінням азоту
при зниженому тиску ($T_{\text{LN2}} = 64,5\,\text{К}$). З кривої інверсії визначають оптимальний тиск
(160 атм), хоча на практиці часто використовують $T = 72\,\text{К}$ і $P = 140\,\text{атм}$ з
технічних міркувань.

\end{document}